\documentclass[11pt]{article}

\usepackage[a4paper,margin=0.85in]{geometry}
\usepackage{amsmath}
\usepackage{booktabs}
\usepackage{array}
\usepackage{xcolor}
\usepackage{hyperref}
\usepackage{fancyhdr}
\usepackage{listings}

\definecolor{navy}{HTML}{18324A}
\definecolor{teal}{HTML}{0B6B72}
\definecolor{green}{HTML}{2E7D32}
\definecolor{red}{HTML}{B3261E}
\definecolor{softgray}{HTML}{F5F7FA}
\definecolor{linegray}{HTML}{D7DEE8}
\definecolor{codebg}{HTML}{F7F9FC}

\hypersetup{
  colorlinks=true,
  linkcolor=teal,
  urlcolor=teal,
  citecolor=teal
}

\pagestyle{fancy}
\fancyhf{}
\lhead{\textcolor{navy}{NNCS-Mamba}}
\rhead{\textcolor{navy}{STL Monitor Report}}
\cfoot{\thepage}
\renewcommand{\headrulewidth}{0.4pt}

\setlength{\parindent}{0pt}
\setlength{\parskip}{6pt}
\setlength{\headheight}{14pt}
\setlength{\tabcolsep}{4pt}
\sloppy

\lstdefinestyle{pythonstyle}{
  backgroundcolor=\color{codebg},
  basicstyle=\ttfamily\small,
  keywordstyle=\color{teal}\bfseries,
  commentstyle=\color{green!60!black},
  stringstyle=\color{red!70!black},
  showstringspaces=false,
  frame=single,
  rulecolor=\color{linegray},
  breaklines=true,
  columns=fullflexible
}
\lstset{style=pythonstyle}

\newcommand{\code}[1]{\texttt{#1}}
\newcommand{\metric}[1]{{\footnotesize\texttt{#1}}}
\newcommand{\rob}{\rho}
\newcommand{\traj}{\xi}
\newcommand{\reportbox}[2]{
  \vspace{4pt}
  \noindent\colorbox{softgray}{
    \parbox{\dimexpr0.98\linewidth-2\fboxsep\relax}{
      \textbf{\textcolor{navy}{#1}}\par
      #2
    }
  }
  \vspace{4pt}
}

\title{
  \vspace{-1.2cm}
  \textbf{\textcolor{navy}{STL Robustness Monitor Implementation Report}}\\
  \large NNCS-Mamba Stabilization Evaluation
}
\author{}
\date{\today}

\begin{document}
\maketitle
\thispagestyle{fancy}

\reportbox{Report Objective}{
This document reports the mechanics of \code{stl\_monitor.py}. The monitor converts one complete drone rollout into quantitative Signal Temporal Logic (STL) robustness scores. The implementation is used to evaluate whether a controller keeps the drone safe, reaches the goal, and remains settled near the goal.
}

\section{Role of the Monitor in the Project}

The STL monitor is not a controller and it does not change the drone action. Its role is evaluation. Given a completed trajectory, it computes a scalar robustness value:

\[
\rob(\traj) > 0 \quad \Rightarrow \quad \text{trajectory satisfies the STL specification}
\]
\[
\rob(\traj) < 0 \quad \Rightarrow \quad \text{trajectory violates at least one part of the specification}
\]

In the training pipeline, a negative robustness value can be treated as a counterexample. Such trajectories are useful for counterexample-guided retraining because they show where the learned controller failed relative to the formal requirement.

\section{Inputs and State Convention}

The main function is:

\begin{lstlisting}[language=Python]
evaluate_stabilization_stl(states, actions, spec, action_low, action_high)
\end{lstlisting}

The expected inputs are summarized below.

\begin{center}
\begin{tabular}{>{\bfseries}p{0.24\linewidth}p{0.66\linewidth}}
\toprule
Input & Meaning \\
\midrule
\code{states} & Array with shape \code{[T, 12]}. It contains the drone state at each timestep of one rollout. \\
\code{actions} & Optional array with shape \code{[T, nu]}. It contains the applied thrust/control input at each timestep. \\
\code{spec} & An \code{STLSpec} object containing the goal position and safety thresholds. \\
\code{action\_low} & Optional lower bounds for each action component. \\
\code{action\_high} & Optional upper bounds for each action component. \\
\bottomrule
\end{tabular}
\end{center}

The monitor uses helper functions from \code{safe\_control\_gym\_config.py} to extract the relevant components:

\[
p_t = [x_t, y_t, z_t], \qquad
v_t = [\dot{x}_t, \dot{y}_t, \dot{z}_t], \qquad
\theta_t = [\phi_t, \theta_t, \psi_t]
\]

Here, \(p_t\) is position, \(v_t\) is velocity, and \(\theta_t\) contains the three attitude angles. The monitor uses all three angles through the maximum absolute angle.

\section{Specification Parameters}

The dataclass \code{STLSpec} stores the numerical requirements used by the monitor.

\begin{center}
\begin{tabular}{>{\bfseries}p{0.36\linewidth}p{0.43\linewidth}p{0.09\linewidth}}
\toprule
Field & Meaning & Default \\
\midrule
\metric{goal\_position} & Target 3D position for stabilization. & required \\
\metric{position\_tolerance} & Radius used to decide whether the drone has reached the goal. & 0.25 \\
\metric{settle\_position\_tolerance} & Radius allowed after the drone has settled. & 0.35 \\
\metric{settle\_speed\_tolerance} & Maximum speed allowed after settling. & 0.25 \\
\metric{x\_bound}, \metric{y\_bound} & Horizontal safety limits. & 2.5 \\
\metric{z\_min}, \metric{z\_max} & Floor and ceiling limits. & 0.0, 2.5 \\
\metric{max\_abs\_angle} & Maximum allowed absolute attitude angle. & 0.75 \\
\metric{reach\_deadline\_step} & Optional deadline for reaching the goal. & None \\
\bottomrule
\end{tabular}
\end{center}

These values define margins, not binary pass/fail labels. A positive margin means the requirement is satisfied with clearance. A negative margin means it is violated.

\section{Implemented STL Formula}

The monitor evaluates the following stabilization requirement:

\[
\varphi =
\mathbf{G}(\text{state safe})
\;\land\;
\mathbf{G}(\text{input safe})
\;\land\;
\mathbf{F}(\text{near goal})
\;\land\;
\mathbf{F}\mathbf{G}(\text{settled})
\]

In words, the trajectory must:

\begin{itemize}
  \item always remain inside the allowed state boundaries,
  \item always respect the input limits when action bounds are provided,
  \item eventually enter the goal region,
  \item eventually become settled and remain settled until the end of the rollout.
\end{itemize}

The final robustness is the minimum robustness among these four parts. This matches the quantitative STL interpretation of conjunction: the weakest requirement determines the robustness of the whole formula.

\section{Temporal Operators}

\subsection{\code{always(signal)}}

\code{always} implements the STL globally operator \(\mathbf{G}\). For a robustness signal \(s_0, s_1, \ldots, s_{T-1}\), it returns:

\[
\mathbf{G}(s) = \min_{0 \leq t < T} s_t
\]

This is used when the requirement must hold for the complete trajectory. One violation is enough to make the global robustness negative.

\begin{lstlisting}[language=Python]
def always(signal: np.ndarray) -> float:
    if len(signal) == 0:
        return -np.inf
    return float(np.min(signal))
\end{lstlisting}

\subsection{\code{eventually(signal, end\_idx=None)}}

\code{eventually} implements the STL finally operator \(\mathbf{F}\). It returns:

\[
\mathbf{F}(s) = \max_{0 \leq t \leq T-1} s_t
\]

If a deadline is supplied, the maximum is only taken up to that timestep:

\[
\mathbf{F}_{[0,d]}(s) = \max_{0 \leq t \leq d} s_t
\]

This is used for reachability: the drone only needs to be close to the goal at least once.

\begin{lstlisting}[language=Python]
def eventually(signal: np.ndarray, end_idx: Optional[int] = None) -> float:
    if len(signal) == 0:
        return -np.inf
    if end_idx is None:
        end_idx = len(signal) - 1
    end_idx = min(end_idx, len(signal) - 1)
    return float(np.max(signal[: end_idx + 1]))
\end{lstlisting}

\subsection{\code{suffix\_always(signal)}}

\code{suffix\_always} computes the robustness of ``always from this timestep to the end'' for every timestep:

\[
\text{suffix\_always}(s)_t = \min_{k \geq t} s_k
\]

This is needed for the formula \(\mathbf{F}\mathbf{G}(\text{settled})\). The monitor first asks: if settling begins at timestep \(t\), does the drone remain settled for the rest of the rollout? Then \code{eventually} checks whether any such timestep exists.

\begin{lstlisting}[language=Python]
def suffix_always(signal: np.ndarray) -> np.ndarray:
    out = np.empty_like(signal, dtype=np.float64)
    running = np.inf
    for idx in range(len(signal) - 1, -1, -1):
        running = min(running, float(signal[idx]))
        out[idx] = running
    return out
\end{lstlisting}

\section{State Safety Robustness}

At each timestep, the monitor computes five safety margins:

\[
m^x_t = x_{\max} - |x_t|
\]
\[
m^y_t = y_{\max} - |y_t|
\]
\[
m^{z,\min}_t = z_t - z_{\min}
\]
\[
m^{z,\max}_t = z_{\max} - z_t
\]
\[
m^\theta_t = \theta_{\max} - \max(|\phi_t|, |\theta_t|, |\psi_t|)
\]

The state safety signal at timestep \(t\) is the smallest of these margins:

\[
s^{state}_t =
\min(m^x_t, m^y_t, m^{z,\min}_t, m^{z,\max}_t, m^\theta_t)
\]

This value represents the closest active state-safety threat at that timestep. The full-trajectory state safety robustness is:

\[
\rob_{state} = \mathbf{G}(s^{state}) = \min_t s^{state}_t
\]

\begin{lstlisting}[language=Python]
state_safety_signal = np.minimum.reduce([
    spec.x_bound - np.abs(positions[:, 0]),
    spec.y_bound - np.abs(positions[:, 1]),
    positions[:, 2] - spec.z_min,
    spec.z_max - positions[:, 2],
    spec.max_abs_angle - np.max(np.abs(angles), axis=1),
])
state_safety = always(state_safety_signal)
\end{lstlisting}

\section{Input Safety Robustness}

Input safety is evaluated only when \code{actions}, \code{action\_low}, and \code{action\_high} are provided. For each action component, two margins are computed:

\[
m^{low}_{t,i} = u_{t,i} - u^{low}_i
\]
\[
m^{high}_{t,i} = u^{high}_i - u_{t,i}
\]

The input safety signal at timestep \(t\) is:

\[
s^{input}_t =
\min_i(m^{low}_{t,i}, m^{high}_{t,i})
\]

The trajectory-level input safety robustness is:

\[
\rob_{input} = \mathbf{G}(s^{input})
\]

If no action bounds are supplied, the implementation sets \(\rob_{input} = \infty\). This means the input-safety term does not constrain the final minimum in that evaluation.

\begin{lstlisting}[language=Python]
low_margin = actions - np.asarray(action_low).reshape(1, -1)
high_margin = np.asarray(action_high).reshape(1, -1) - actions
input_safety = always(np.minimum(
    np.min(low_margin, axis=1),
    np.min(high_margin, axis=1),
))
\end{lstlisting}

\section{Reach Robustness}

The distance to the goal is:

\[
d_t = \|p_t - p_{goal}\|_2
\]

The reach signal is:

\[
s^{reach}_t = r_{reach} - d_t
\]

where \(r_{reach}\) is \code{position\_tolerance}. The reach requirement is satisfied if this signal becomes positive at least once:

\[
\rob_{reach} = \mathbf{F}(s^{reach})
\]

If a deadline is configured, the maximum is only taken before or at that deadline.

\begin{lstlisting}[language=Python]
distance_to_goal = np.linalg.norm(
    positions - spec.goal_position.reshape(1, 3),
    axis=1,
)
reach_signal = spec.position_tolerance - distance_to_goal
eventually_reach = eventually(reach_signal, spec.reach_deadline_step)
\end{lstlisting}

\section{Settled Robustness}

The monitor separates reaching the goal from stabilizing at the goal. A drone may briefly pass through the target region but still fail to hover there. For this reason, settling uses both position and speed:

\[
s^{settle}_t =
\min(r_{settle} - d_t,\; v_{settle} - \|v_t\|_2)
\]

The drone is considered settled at timestep \(t\) only when it is close enough to the goal and slow enough.

The full settled requirement is:

\[
\rob_{settled} =
\mathbf{F}\mathbf{G}(s^{settle})
\]

This means that there must exist a timestep after which the drone remains settled until the rollout ends.

\begin{lstlisting}[language=Python]
settled_signal = np.minimum(
    spec.settle_position_tolerance - distance_to_goal,
    spec.settle_speed_tolerance - speed,
)
eventually_always_settled = eventually(suffix_always(settled_signal))
\end{lstlisting}

\section{Final Robustness Calculation}

The final robustness value is:

\[
\rob =
\min(
\rob_{state},
\rob_{input},
\rob_{reach},
\rob_{settled}
)
\]

The implementation returns both the final score and the component scores:

\begin{lstlisting}[language=Python]
robustness = min(
    state_safety,
    input_safety,
    eventually_reach,
    eventually_always_settled,
)
\end{lstlisting}

\begin{center}
\begin{tabular}{>{\bfseries}p{0.47\linewidth}p{0.43\linewidth}}
\toprule
Returned key & Interpretation \\
\midrule
\metric{stl\_robustness} & Final robustness of the full STL formula. \\
\metric{stl\_satisfied} & 1.0 if final robustness is non-negative, otherwise 0.0. \\
\metric{state\_safety\_robustness} & Worst state-safety margin over the trajectory. \\
\metric{input\_safety\_robustness} & Worst input-limit margin over the trajectory. \\
\metric{eventually\_reach\_robustness} & Best goal-reaching margin over the trajectory. \\
\metric{eventually\_always\_settled\_robustness} & Best timestep from which the drone remains settled. \\
\metric{min\_distance\_to\_goal} & Closest distance to the goal during the rollout. \\
\metric{final\_distance\_to\_goal} & Distance to the goal at the final timestep. \\
\bottomrule
\end{tabular}
\end{center}

\section{Function-by-Function Summary}

\begin{center}
\begin{tabular}{>{\bfseries}p{0.39\linewidth}p{0.51\linewidth}}
\toprule
Function or object & Purpose in the monitor \\
\midrule
\code{STLSpec} & Stores the formal thresholds for the stabilization task. \\
\code{always} & Implements global STL robustness by taking the minimum over time. \\
\code{eventually} & Implements reachability STL robustness by taking the maximum over time. \\
\code{suffix\_always} & Computes future-global robustness from every timestep, used for settling. \\
\metric{evaluate\_stabilization\_stl} & Computes state safety, input safety, reach, settled behavior, and final robustness. \\
\bottomrule
\end{tabular}
\end{center}

\section{Counterexample Interpretation}

For CEGIS-style retraining, the important decision is:

\[
\rob(\traj) < 0
\]

If the final robustness is negative, then at least one required part of the STL formula failed. The complete rollout can be stored as a counterexample because it demonstrates a concrete failure mode of the current controller. The component robustness values identify the cause of failure, for example:

\begin{itemize}
  \item negative state safety: boundary, floor, ceiling, or angle violation;
  \item negative input safety: action limit violation;
  \item negative reach robustness: goal region was never reached;
  \item negative settled robustness: the drone did not remain close and slow after reaching the goal.
\end{itemize}

\section{Validation Notes and Limitations}

The implementation is intentionally small and inspectable. It uses standard quantitative STL operations: minimum for conjunction/global requirements and maximum for eventually requirements.

Current limitations are:

\begin{itemize}
  \item The monitor is specialized to stabilization-style tasks around one goal position.
  \item The current formula does not yet encode waypoint sequences or docking geometry.
  \item The monitor assumes the state ordering supplied by \code{safe\_control\_gym\_config.py}; any environment change must preserve or update that mapping.
  \item The returned robustness is only as meaningful as the chosen thresholds in \code{STLSpec}.
\end{itemize}

The main validation requirement is to test the formulas on hand-constructed trajectories: clearly safe trajectories should produce positive robustness, and trajectories with known violations should produce negative robustness in the corresponding component.

\end{document}
